\subsection{N-gram Embedding}
\label{sec:ngram-emb}
\paragraph{Motivation.}
Embedding-based memory offers a complementary dimension for scaling model capacity~\citep{Jordan2024,google2025gemma3n,rwkvcommunity2025deepembed,gemmateam2026gemma4,sadhukhan2026stem,tseng2026l3}.
N-gram embeddings further generalize unigram lookup by conditioning memory retrieval on local context rather than token identity alone~\citep{huang2021lookuplm,roy2022ngrammer,huang2025overtokenized,yu2025scone,cheng2026engram, liu2026scalingembeddings,chen2026xgram}. 
Concretely, short n-grams ending at each token serve as keys into embedding tables, and the retrieved vectors augment the corresponding token representation.
N-gram memory scales capacity with negligible additional per-token FLOPs, while deterministic addressing enables host-memory offloading and asynchronous prefetching.
In this section, we systematically study the key architectural choices for N-gram embedding.
We use multi-head hashing for embedding lookup~\citep{roy2022ngrammer,huang2025overtokenized}. 
The embeddings are further injected into residual stream through contextual gating mechanism~\citep{cheng2026engram}.
We use 300 tokens per active parameter (TPP) throughout experiments.

\subsubsection{Placement of the N-gram Embedding Layers}
% \begin{table*}[t]
% \centering
% \caption{Effect of N-gram embedding layer placement. The total number of N-gram embedding parameters
% is fixed across all settings.}
% \label{tab:ngram-layer-index}
% \resizebox{\textwidth}{!}{
% \begin{tabular}{l c cccc ccc c c cc c}
% \toprule
% \multirow{2}{*}{Layer Index}
% & \multirow{2}{*}{Loss}
% & \multicolumn{4}{c}{Knowledge}
% & \multicolumn{3}{c}{STEM}
% & \multicolumn{1}{c}{Reasoning}
% & \multicolumn{1}{c}{Multilingual}
% & \multicolumn{2}{c}{Code}
% & \multirow{2}{*}{Avg.} \\
% \cmidrule(lr){3-6}
% \cmidrule(lr){7-9}
% \cmidrule(lr){10-10}
% \cmidrule(lr){11-11}
% \cmidrule(lr){12-13}
% &
% & MMLU
% & MMLU-Redux
% & MMLU-Pro
% & SuperGPQA
% & MATH
% & GSM8K
% & GPQA-Diamond
% & BBH
% & MMMLU
% & EvalPlus
% & MultiPL-E
% & \\
% \midrule
% w/o N-gram emb. & 1.585 & 62.78 & 60.62 & 33.43 & 20.97 & 32.52 & 59.21 & 25.34 & 53.40 & 54.06 & 52.13 & 45.42 & 45.44 \\
% 1st & 1.541 & 64.19 & 62.17 & 35.25 & 21.30 & 36.20 & \bf{65.73} & 24.91 & 56.00 & 56.16 & 50.95 & \bf{47.45} & 47.30 \\
% 2nd & 1.541 & 64.71 & 63.28 & \bf{35.80} & 21.49 & 37.32 & 64.00 & 27.89 & 57.56 & \bf{56.64} & \bf{53.09} & 45.56 & 47.94 \\
% 3rd & 1.543 & 63.20 & 62.02 & 34.93 & 20.67 & 35.74 & 63.15 & 27.84 & 56.71 & 55.02 & 50.15 & 44.90 & 46.76 \\
% 4th & 1.544 & 63.36 & 61.48 & 34.81 & \bf{22.33} & 36.20 & 61.26 & 28.63 & 55.32 & 55.79 & 51.05 & 45.60 & 46.89 \\
% 10th & 1.544 & 64.22 & 62.38 & 34.99 & 20.81 & 33.80 & 62.17 & 27.75 & 55.54 & 54.81 & 52.69 & 43.61 & 46.62 \\
% 15th & 1.543 & \bf{65.07} & \bf{63.57} & 34.95 & 21.35 & 36.12 & 63.50 & 26.80 & 57.15 & 55.98 & 52.21 & 44.42 & 47.37 \\
% 25th & 1.541 & 64.70 & 62.24 & 35.15 & 22.13 & 36.26 & 63.31 & \bf{28.85} & 55.73 & 55.99 & 52.66 & 44.33 & 47.40 \\
% \midrule
% 2nd + 15th & 1.541 & 63.82 & 61.16 & 35.63 & 21.25 & \bf{37.48} & 63.23 & 28.00 & 56.52 & 55.68 & 50.44 & 43.85 & 47.01 \\
% 2nd + 25th & \bf{1.540} & 64.94 & 62.49 & 35.40 & 21.80 & 37.40 & 64.33 & 28.28 & \bf{57.79} & 56.45 & 51.69 & 44.73 & 47.75 \\
% \bottomrule
% \end{tabular}
% }
% \end{table*}
\begin{table*}[h]
\centering
\caption{Effect of N-gram embedding layer placement. The total number of N-gram embedding parameters
is fixed across all settings.}
\label{tab:ngram-layer-index}
\resizebox{\textwidth}{!}{
\begin{tabular}{l c ccc cc c c cc c}
\toprule
&
& \multicolumn{3}{c}{Knowledge}
& \multicolumn{2}{c}{STEM}
& \multicolumn{1}{c}{Reasoning}
& \multicolumn{1}{c}{Multilingual}
& \multicolumn{2}{c}{Code}
&  \\
\cmidrule(lr){3-5}
\cmidrule(lr){6-7}
\cmidrule(lr){8-8}
\cmidrule(lr){9-9}
\cmidrule(lr){10-11}

Layer Index & Loss
& MMLU
& MMLU-Pro
& SuperGPQA
& MATH
& GSM8K
& BBH
& MMMLU
& EvalPlus
& MultiPL-E
& Avg. \\
\midrule
w/o N-gram emb.
& 1.585 & 62.78 & 33.43 & 20.97 & 32.52 & 59.21 & 53.40 & 54.06 & 52.13 & 45.42 & 45.44 \\
1st
& 1.541 & 64.19 & 35.25 & 21.30 & 36.20 & \bf{65.73} & 56.00 & 56.16 & 50.95 & \bf{47.45} & 47.30 \\
2nd
& 1.541 & 64.71 & \bf{35.80} & 21.49 & 37.32 & 64.00 & 57.56 & \bf{56.64} & \bf{53.09} & 45.56 & 47.94 \\
3rd
& 1.543 & 63.20 & 34.93 & 20.67 & 35.74 & 63.15 & 56.71 & 55.02 & 50.15 & 44.90 & 46.76 \\
4th
& 1.544 & 63.36 & 34.81 & \bf{22.33} & 36.20 & 61.26 & 55.32 & 55.79 & 51.05 & 45.60 & 46.89 \\
10th
& 1.544 & 64.22 & 34.99 & 20.81 & 33.80 & 62.17 & 55.54 & 54.81 & 52.69 & 43.61 & 46.62 \\
15th
& 1.543 & \bf{65.07} & 34.95 & 21.35 & 36.12 & 63.50 & 57.15 & 55.98 & 52.21 & 44.42 & 47.37 \\
25th
& 1.541 & 64.70 & 35.15 & 22.13 & 36.26 & 63.31 & 55.73 & 55.99 & 52.66 & 44.33 & 47.40 \\
\midrule
2nd + 15th
& 1.541 & 63.82 & 35.63 & 21.25 & \bf{37.48} & 63.23 & 56.52 & 55.68 & 50.44 & 43.85 & 47.01 \\
2nd + 25th
& \bf{1.540} & 64.94 & 35.40 & 21.80 & 37.40 & 64.33 & \bf{57.79} & 56.45 & 51.69 & 44.73 & 47.75 \\
\bottomrule
\end{tabular}
}
\end{table*}
We ablate the placement and number of N-gram embedding layers under a fixed parameter budget.
Single-layer variants span shallow (Layers 1-4), intermediate (Layers 10 and 15), and deep (Layer 25) locations.
For multi-layer configurations, we combine a shallow layer (Layer 2) with either an intermediate layer (Layer 15) or a deep layer (Layer 25).
The results are shown in Table~\ref{tab:ngram-layer-index}.

No single depth regime consistently dominates. 
The first two layers perform strongly, while intermediate and deep placements remain competitive.
Distributing the same parameter budget across multiple layers yields no consistent benefit. 
The marginal loss reduction from combining Layers 2 and 25 does not translate into improved downstream performance. 
A single N-gram embedding layer is sufficient. 
Moreover, the relative performance of different placements is similar under full attention and GDN (\S\ref{sec:gdn-hybrid}), suggesting that placement choice is largely insensitive to the attention mechanism.
We place it at Layer 2, \textit{allowing host-memory prefetching to overlap with the computation of the first layer.}

\subsubsection{N-gram Vocabulary Size}
We further explore the effect of scaling the vocabulary size of N-gram embeddings.

% \begin{table*}[t]
% \centering
% \caption{Effect of N-gram vocabulary scaling under a fixed total model parameter budget. Vocabulary scales are measured relative to the base tokenizer vocabulary size (250K). The number of MoE experts is adjusted to offset the additional N-gram embedding parameters.}
% \label{tab:ngram-vocab-size-samebudget}
% \resizebox{\textwidth}{!}{
% \begin{tabular}{l *{13}{c}}
% \toprule
% Vocab. Scale (Param. Ratio)
% & Loss & Uncheatable PPL 
% & MMLU
% & MMLU-Redux
% & MMLU-Pro
% & SuperGPQA
% & MATH
% & GSM8K
% & GPQA-Diamond
% & BBH
% & C-Eval
% & CMMLU
% & MMMLU\\
% \midrule
% None (0\%) & 1.202 & 5.55 & \bf{68.25} & \bf{67.28} & 44.38 & 25.22 & 46.02 & \bf{74.79} & \bf{30.74} & \bf{65.71} & 70.78 & 73.01 & \bf{58.32} \\
% 5$\times$ (10\%) & 1.200 & \bf{5.54} & 68.15 & 66.20 & 44.49 & 24.25 & 45.64 & 72.86 & 29.99 & 65.39 & 70.93 & \bf{73.44} & 57.49 \\
% 10$\times$ (25\%) & \bf{1.197} & 5.55 & 67.71 & 66.88 & \bf{44.66} & \bf{25.64} & \bf{46.56} & 73.65 & 30.15 & 65.54 & 70.71 & 73.31 & 56.22 \\
% 30$\times$ (50\%) & 1.201 & 5.59 & 67.75 & 65.66 & 42.61 & 24.18 & 44.62 & 74.45 & 29.29 & 65.55 & \bf{72.49} & 73.28 & 56.66 \\
% \bottomrule
% \end{tabular}
% }
% \end{table*}
\begin{table*}[h]
\centering
\caption{Effect of N-gram vocabulary scaling under a fixed total model parameter budget. Vocabulary scales are measured relative to the base tokenizer vocabulary size (250K). The number of MoE experts is adjusted to offset the additional N-gram embedding parameters.}
\label{tab:ngram-vocab-size-samebudget}
\resizebox{\textwidth}{!}{
\begin{tabular}{l cc ccc cc c cc c}
\toprule
Vocab. Scale
& \multirow{2}{*}{Loss}
& Uncheatable
& \multicolumn{3}{c}{Knowledge}
& \multicolumn{2}{c}{STEM}
& \multicolumn{1}{c}{Reasoning}
& \multicolumn{2}{c}{Chinese}
& \multicolumn{1}{c}{Multilingual} \\
\cmidrule(lr){4-6}
\cmidrule(lr){7-8}
\cmidrule(lr){9-9}
\cmidrule(lr){10-11}
\cmidrule(lr){12-12}
(Param. Ratio) & & PPL & MMLU
& MMLU-Pro
& SuperGPQA
& MATH
& GSM8K
& BBH
& C-Eval
& CMMLU
& MMMLU \\
\midrule
None (0\%)
& 1.202 & 5.55 & \bf{68.25} & 44.38 & 25.22 & 46.02 & \bf{74.79} & \bf{65.71} & 70.78 & 73.01 & \bf{58.32} \\
5$\times$ (10\%)
& 1.200 & \bf{5.54} & 68.15 & 44.49 & 24.25 & 45.64 & 72.86 & 65.39 & 70.93 & \bf{73.44} & 57.49 \\
10$\times$ (25\%)
& \bf{1.197} & 5.55 & 67.71 & \bf{44.66} & \bf{25.64} & \bf{46.56} & 73.65 & 65.54 & 70.71 & 73.31 & 56.22 \\
30$\times$ (50\%)
& 1.201 & 5.59 & 67.75 & 42.61 & 24.18 & 44.62 & 74.45 & 65.55 & \bf{72.49} & 73.28 & 56.66 \\
\bottomrule
\end{tabular}
}
\end{table*}
\paragraph{Allocation under Fixed Parameter Budget.}
Following prior work, we scale the N-gram embedding slots while reducing the number of experts to maintain a fixed total parameter budget. 
The results are reported in Table~\ref{tab:ngram-vocab-size-samebudget}. 
The loss varies non-monotonically with vocabulary size and is lowest at $10\times$ ($25\%$), consistent with the allocation sweet spots reported in prior work~\citep{liu2026scalingembeddings,cheng2026engram}.
The same optimum is not evident in other evaluations. The out-of-domain uncheatable PPL changes little across budgets, and downstream benchmarks show no clear improvement over the MoE-only baseline.
These results suggest that \textit{N-gram embeddings and MoE experts play distinct roles in scaling capacity.
We therefore study N-gram vocabulary scaling while holding the MoE parameter budget fixed.}

\paragraph{Scaling with Additional Parameters.}
% \begin{table*}[t]
% \centering
% \caption{Effect of N-gram vocabulary scaling. Vocabulary scales are measured relative to the base tokenizer vocabulary size (250K). Larger vocabularies increase the total number of parameters.}
% \label{tab:ngram-vocab-size}
% \resizebox{\textwidth}{!}{
% \begin{tabular}{l *{12}{c}}
% \toprule
% Vocab. Scale
% & Loss
% & MMLU
% & MMLU-Redux
% & MMLU-Pro
% & SuperGPQA
% & MATH
% & GSM8K
% & GPQA-Diamond
% & BBH
% & C-Eval
% & CMMLU
% & MMMLU\\
% \midrule
% None & 1.585 & 62.78 & 60.62 & 33.43 & 20.97 & 32.52 & 59.21 & 25.34 & 53.40 & 66.91 & 68.10 & 54.06 \\
% 20$\times$ & 1.553 & 64.14 & \bf{63.32} & 34.46 & 21.83 & \bf{37.38} & \bf{65.09} & \bf{29.14} & 57.13 & 71.75 & 72.29 & 55.94 \\
% 50$\times$ & 1.541 & 64.71 & 63.28 & 35.80 & 21.49 & 37.32 & 64.00 & 27.89 & \bf{57.56} & 72.12 & 72.48 & 56.64 \\
% 100$\times$ & 1.534 & 64.70 & 62.17 & \bf{35.87} & \bf{21.93} & 36.98 & 63.08 & 29.04 & 56.03 & 73.75 & 72.73 & \bf{56.65} \\
% 200$\times$ & \bf{1.526} & \bf{64.85} & 62.56 & 35.21 & 21.11 & 35.34 & 62.96 & 27.15 & 56.23 & \bf{74.94} & \bf{73.24} & 55.82 \\
% \bottomrule
% \end{tabular}
% }
% \end{table*}

\begin{table*}[h]
\centering
\caption{Effect of N-gram vocabulary scaling. Vocabulary scales are measured relative to the base tokenizer vocabulary size (250K). Larger vocabularies increase the total number of parameters.}
\label{tab:ngram-vocab-size}
\resizebox{\textwidth}{!}{
\begin{tabular}{l c ccc cc c cc c}
\toprule
& 
& \multicolumn{3}{c}{Knowledge}
& \multicolumn{2}{c}{STEM}
& \multicolumn{1}{c}{Reasoning}
& \multicolumn{2}{c}{Chinese}
& \multicolumn{1}{c}{Multilingual} \\
\cmidrule(lr){3-5}
\cmidrule(lr){6-7}
\cmidrule(lr){8-8}
\cmidrule(lr){9-10}
\cmidrule(lr){11-11}

Vocab. Scale & Loss & MMLU
& MMLU-Pro
& SuperGPQA
& MATH
& GSM8K
& BBH
& C-Eval
& CMMLU
& MMMLU \\
\midrule
None
& 1.585 & 62.78 & 33.43 & 20.97 & 32.52 & 59.21 & 53.40 & 66.91 & 68.10 & 54.06 \\
20$\times$
& 1.553 & 64.14 & 34.46 & 21.83 & \bf{37.38} & \bf{65.09} & 57.13 & 71.75 & 72.29 & 55.94 \\
50$\times$
& 1.541 & 64.71 & 35.80 & 21.49 & 37.32 & 64.00 & \bf{57.56} & 72.12 & 72.48 & 56.64 \\
100$\times$
& 1.534 & 64.70 & \bf{35.87} & \bf{21.93} & 36.98 & 63.08 & 56.03 & 73.75 & 72.73 & \bf{56.65} \\
200$\times$
& \bf{1.526} & \bf{64.85} & 35.21 & 21.11 & 35.34 & 62.96 & 56.23 & \bf{74.94} & \bf{73.24} & 55.82 \\
\bottomrule
\end{tabular}
}
\end{table*}

Because embedding tables are sparsely accessed and deterministically addressed, they can be scaled with negligible additional per-token computation and stored in off-accelerator storage~\citep{google2025gemma3n,cheng2026engram}.
We therefore move beyond the fixed-size setting and scale the N-gram vocabulary from $20V$ to $200V$, where $V$ denotes the base tokenizer vocabulary size of Qwen3.5~\citep{qwenteam2026qwen35}.
The results are shown in Table~\ref{tab:ngram-vocab-size}.
Loss decreases monotonically as the N-gram vocabulary grows, while downstream performance does not follow the same trend.
Vocabulary scaling yields broad gains over the baseline, but performance on some benchmarks saturates or fluctuates as the vocabulary grows.
Furthermore, performance on Chinese benchmarks (e.g., C-Eval and CMMLU) improves consistently with N-gram vocabulary size.

Beyond vocabulary scaling, we explored a range of strategies for improving parameter efficiency of N-gram embeddings, including but not limited to token normalization for vocabulary compression~\citep{cheng2026engram}, non-uniform allocation across N-gram orders, and frequency-based partitioning of embedding slots.
Despite these efforts, we observed no consistent performance gains in our training recipe.
